% =========================================================
%  sesi/02/sesi.tex
%  Sesi 2 — Level Rak: Tilemap, Physics, dan Kontrol Pemain
% =========================================================

\sesi{02}{Level Rak: Tilemap, Physics, dan Kontrol Pemain}{3.2}{140 menit}

Ini pertemuan minggu kedua belas. Scene \file{LevelRak} yang Anda susun pada
sesi 1 masih berupa balok-balok primitif: belum ada tabrakan, belum ada
gravitasi, dan pemain belum dapat digerakkan. Sesi ini mengubahnya menjadi
level yang benar-benar dapat dimainkan.

Pekerjaan pertama adalah menyusun lantai dan rak dengan \istilah{Tilemap},
yaitu lembaran petak yang diisi pola. Pekerjaan kedua memberikan sifat fisika:
lantai dan rak menjadi benda diam yang menghalangi, sedangkan pemain menjadi
benda bergerak yang tunduk pada gravitasi. Pekerjaan ketiga menulis skrip
kontrol pemain, lalu menambahkan artefak dan pintu sebagai daerah
\istilah{trigger}.

Ketiga pekerjaan itu mengikuti Bab XIII. Urutannya sengaja dipertahankan:
menyusun ruang, memberi bentuk batas, baru menulis perilaku. Apabila urutan
itu dibalik, Anda akan menghabiskan waktu mencari kesalahan kode padahal
penyebabnya berada pada komponen yang belum dipasang.

% ---------------------------------------------------------
\bagian{Tujuan Sesi}

\begin{tujuanbox}
Setelah menuntaskan sesi ini, Anda dapat menunjukkan:
\begin{enumerate}[leftmargin=1.5em, itemsep=0.2em]
  \item scene \file{LevelRak} yang lantai dan raknya tersusun dari
        \istilah{Tilemap} dan dapat dipijak;
  \item pemain berkomponen \istilah{Rigidbody2D} dan \istilah{Capsule
        Collider2D} yang berjalan, melompat, dan mendarat tanpa menembus lantai;
  \item skrip \code{KontrolPemain} yang terpasang dengan seluruh rujukan
        \menu{Inspector} terisi;
  \item tiga artefak yang lenyap ketika disentuh dan pintu yang melaporkan
        keadaannya pada \menu{Console};
  \item penjelasan mengapa masukan dibaca pada \code{Update} sedangkan gerak
        fisika dikerjakan pada \code{FixedUpdate}.
\end{enumerate}
\end{tujuanbox}

% ---------------------------------------------------------
\bagian{Prasyarat}

\begin{prasyaratbox}
\begin{itemize}[leftmargin=1.4em, itemsep=0.25em]
  \item Checkpoint sesi 1 sudah terbuka: project \NamaProject\ dengan scene
        \file{LevelRak} berisi \code{Gudang}, lantai, tiga rak, pemain, dan
        pintu.
  \item Anda telah membaca Bab XIII bagian 13.1 sampai 13.5, yaitu bagian
        tentang penyusunan level, \istilah{Collider} dan \istilah{Rigidbody},
        kontrol pemain, serta \istilah{trigger}.
  \item Scene sesi 1 tersimpan. Perubahan pada sesi ini akan menimpa scene
        itu, sehingga simpanlah salinan bila Anda masih memerlukannya.
\end{itemize}
\end{prasyaratbox}

% ---------------------------------------------------------
\bagian{Peta Kemampuan}

\petasesi
  {Scene primitif tanpa tabrakan dan tanpa kontrol. Pemain hanya benda
   gambar yang diam.}
  {\istilah{Tilemap}, \istilah{Collider2D}, \istilah{Rigidbody2D},
   \istilah{trigger}, dan skrip kontrol pemain.}
  {Penggantian \istilah{greybox} dengan aset, \istilah{Animator}, dan
   \istilah{Canvas} antarmuka pada sesi 3.}

% ---------------------------------------------------------
\bagian{Konsep yang Dipakai}

\subsection*{Tilemap: satu lembar, satu bentuk batas}

\istilah{Tilemap} adalah lembaran petak yang dipakai untuk menyusun lantai,
dinding, dan rak. Keuntungannya bukan pada tampilan, melainkan pada bentuk
batasnya: seluruh petak yang tersusun cukup memakai satu \istilah{Collider}.
Komponen \code{Tilemap Collider 2D} membentuk batas menurut petak yang
dilukis, lalu \code{Composite Collider 2D} menggabungkan batas-batas kecil itu
menjadi satu bentuk utuh. Tanpa penggabungan, pemain dapat tersangkut pada
garis sambung antarpetak.

\subsection*{Benda diam dan benda bergerak}

Lantai, rak, dan pintu tidak bergerak. Mereka cukup memiliki \istilah{Collider}
tanpa \code{Rigidbody2D}. Pemain bergerak, sehingga memerlukan keduanya:
\istilah{Rigidbody2D} membuatnya mengikuti gravitasi dan \istilah{Collider2D}
memberinya bentuk batas. Memberi \code{Rigidbody2D} kepada benda diam hanya
menambah perhitungan tanpa manfaat.

\subsection*{Collider biasa dan trigger}

\istilah{Collider} biasa menghalangi gerak. \istilah{Trigger} tidak
menghalangi apa pun; ia hanya melaporkan bahwa ada benda lain yang masuk.
Artefak dan pintu memakai trigger supaya pemain dapat melewatinya sambil
memicu sesuatu. Tiga syarat agar laporannya sampai: pilihan \menu{Is Trigger}
dicentang, sekurang-kurangnya satu benda memiliki \code{Rigidbody2D}, dan
kedua benda berada pada \istilah{layer} yang saling diperiksa.

\subsection*{Dua fungsi, dua pekerjaan}

\code{Update} dipanggil setiap bingkai dengan selang yang berubah-ubah, dan
dipakai untuk membaca masukan. \code{FixedUpdate} dipanggil dengan selang
tetap, dan dipakai untuk menggerakkan \code{Rigidbody2D} serta memeriksa tanah.
Memindahkan \code{Rigidbody2D} melalui \code{transform.position} pada
\code{Update} membuat benda bergetar dan menembus dinding, sebab perubahan itu
tidak diperhitungkan oleh langkah fisika.

Pemeriksaan tanah tidak membandingkan posisi dengan letak lantai. Ia memeriksa
lingkaran kecil di bawah kaki pemain melalui \code{Physics2D.OverlapCircle}.
Cara itu tetap bekerja pada lantai bertingkat maupun permukaan miring.
Lingkaran pemeriksaannya digambar pada tampilan \menu{Scene} agar mudah
diperiksa.

% ---------------------------------------------------------
\bagian{Tutorial}

Nama menu dan jendela ditulis seperti yang tampak pada Unity LTS. Kerjakan
berurutan; jangan melompat ke bagian skrip sebelum bentuk batas level selesai.

\begin{langkah}
  \item \textbf{Buka scene lama.} Pada jendela \menu{Project}, buka
        \file{Assets/Scenes}, lalu klik dua kali \file{LevelRak}. Pastikan
        Hierarchy masih memuat \code{Gudang} beserta anak-anaknya.

  \item \textbf{Buat Grid dan tilemap lantai.} Klik kanan pada
        \menu{Hierarchy}, pilih \menu{2D Object > Tilemap > Rectangular}. Beri
        nama \code{Grid} pada objek induk dan \code{TilemapLantai} pada
        anaknya. Pada \menu{Inspector} Grid, setel \menu{Cell Size} menjadi
        $(1,1,0)$.

  \item \textbf{Buat tilemap rak.} Klik kanan \code{Grid}, pilih
        \menu{2D Object > Tilemap > Rectangular}, lalu beri nama
        \code{TilemapRak}. Kedua tilemap menjadi anak \code{Grid} agar
        petaknya sejajar.

  \item \textbf{Buat aset petak.} Pada jendela \menu{Project}, klik kanan
        folder \file{Assets}, pilih \menu{Create > Folder}, dan beri nama
        \code{Tiles}. Klik kanan folder itu, pilih
        \menu{Create > 2D > Tiles > Tile}, lalu beri nama \code{TileLantai}.
        Pada \menu{Inspector}, isi \menu{Sprite} dengan sprite \code{Square}
        dan \menu{Color} dengan \code{#46587E}. Gandakan dengan
        \menu{Ctrl+D}, beri nama \code{TileRak}, lalu ubah warnanya menjadi
        \code{#60749E}.

  \item \textbf{Buka palet.} Pilih menu \menu{Window > 2D > Tile Palette}.
        Tekan \menu{Create New Palette}, beri nama \code{PaletGudang}, dan
        simpan di folder \file{Assets/Tiles}. Seret \code{TileLantai} dan
        \code{TileRak} dari jendela \menu{Project} ke dua sel pertama palet.

  \item \textbf{Lukis lantai.} Pilih \code{TilemapLantai} pada
        \menu{Hierarchy}. Pada jendela \menu{Tile Palette}, pilih
        \code{TileLantai} dan alat \menu{Brush} (tombol \menu{B}). Pada
        tampilan \menu{Scene}, lukis satu baris petak pada $y=-4$ dari $x=-9$
        sampai $x=9$.

  \item \textbf{Lukis tiga rak.} Pilih \code{TilemapRak}, lalu pilih
        \code{TileRak} pada palet. Lukis tiga rak masing-masing selebar empat
        petak: $x=-6$ sampai $-3$ pada $y=-2$; $x=-1$ sampai $2$ pada $y=0$;
        dan $x=4$ sampai $7$ pada $y=-2$.

  \item \textbf{Buat layer Ground.} Pilih menu
        \menu{Edit > Project Settings > Tags and Layers}. Pada bagian
        \menu{Layers}, isi \menu{User Layer 6} dengan \code{Ground}. Kembali
        ke \menu{Hierarchy}, setel \menu{Layer} kedua tilemap menjadi
        \code{Ground}.

  \item \textbf{Pasang collider tilemap.} Pilih \code{TilemapLantai}, tekan
        \menu{Add Component}, tambahkan \menu{Tilemap Collider 2D}, lalu
        tambahkan \menu{Composite Collider 2D}. Komponen \code{Rigidbody2D}
        ikut terpasang; setel \menu{Body Type} menjadi \menu{Static} dan
        \menu{Geometry Type} menjadi \menu{Polygons}. Pada
        \code{Tilemap Collider 2D}, centang \menu{Used By Composite}. Ulangi
        langkah yang sama untuk \code{TilemapRak}.

  \item \textbf{Berikan fisika kepada pemain.} Pilih \code{Pemain}, tambahkan
        \menu{Rigidbody2D} dengan \menu{Gravity Scale} 2.5,
        \menu{Collision Detection} \menu{Continuous}, dan
        \menu{Interpolate} \menu{Interpolate}. Pada \menu{Constraints},
        centang \menu{Freeze Rotation Z}. Tambahkan \menu{Capsule Collider2D}
        dengan \menu{Direction} \menu{Vertical} dan \menu{Size} $(0.7,0.9)$.

  \item \textbf{Beri tag Player.} Pada \menu{Inspector} pemain, buka daftar
        \menu{Tag}, pilih \menu{Add Tag}, lalu buat tag \code{Player}. Kembali
        ke pemain dan setel tag-nya menjadi \code{Player}.

  \item \textbf{Buat titik pemeriksaan tanah.} Klik kanan \code{Pemain},
        pilih \menu{Create Empty}, beri nama \code{TitikTanah}, dan setel
        \menu{Position} lokalnya menjadi $(0,-0.5,0)$.
\end{langkah}

\begin{ujibox}
Sebelum menulis skrip, tekan \menu{Play}. Periksa empat hal berikut, lalu
hentikan \menu{Play}.
\begin{itemize}[leftmargin=1.4em, itemsep=0.2em]
  \item Lantai dan tiga rak tampak sebagai deretan petak, bukan balok
        primitif.
  \item Pemain jatuh dan berhenti di atas lantai, tidak menembusnya.
  \item Pemain dapat berdiri di atas rak tanpa tergelincir turun.
  \item Pada menu \menu{Window > Analysis > Physics Debug}, bentuk batas
        lantai dan rak tampak sebagai satu bentuk utuh, bukan petak-petak
        terpisah.
\end{itemize}
\end{ujibox}

\begin{langkah}
  \item \textbf{Buat tiga skrip.} Pada folder \file{Assets/Scripts}, klik
        kanan lalu pilih \menu{Create > C\# Script} tiga kali dengan nama
        \code{KontrolPemain}, \code{Artefak}, dan \code{Pintu}. Buka setiap
        berkas dan salin isinya dari bagian Kode Sesi Ini. Nama kelas harus
        sama dengan nama berkas.

  \item \textbf{Pasang kontrol pemain.} Seret \code{KontrolPemain} dari
        jendela \menu{Project} ke objek \code{Pemain}. Pada
        \menu{Inspector}, isi \menu{Ground Check} dengan \code{TitikTanah},
        \menu{Ground Layer} dengan \code{Ground}, \menu{Move Speed} 6,
        \menu{Jump Speed} 11, dan \menu{Ground Radius} 0.15.

  \item \textbf{Buat tiga artefak.} Klik kanan \code{Gudang}, pilih
        \menu{2D Object > Sprites > Square}, lalu beri nama
        \code{Artefak1}. Setel \menu{Position} $(-4,-1.1,0)$, \menu{Scale}
        $(0.5,0.5,1)$, dan \menu{Color} \code{#DEB242}. Tambahkan
        \menu{Circle Collider2D} dengan \menu{Radius} 0.25 dan centang
        \menu{Is Trigger}, lalu tambahkan komponen \code{Artefak}. Ulangi
        untuk \code{Artefak2} pada $(1,0.9,0)$ dan \code{Artefak3} pada
        $(6,-1.1,0)$.

  \item \textbf{Ubah pintu menjadi trigger.} Pilih \code{Pintu}, setel
        \menu{Position} menjadi $(8.3,-2.6,0)$, lalu tambahkan
        \menu{Box Collider2D} dengan \menu{Size} $(1,1.8)$ dan centang
        \menu{Is Trigger}. Tambahkan komponen \code{Pintu} dan isi
        \menu{Artefak Diperlukan} dengan 3.

  \item \textbf{Uji levelnya.} Tekan \menu{Play}. Gerakkan pemain dengan
        tombol panah atau \menu{A}/\menu{D}, lompat dengan \menu{Space},
        sentuh ketiga artefak, lalu masuk ke daerah pintu. Amati
        \menu{Console}, lalu simpan scene dengan \menu{Ctrl+S}.
\end{langkah}

\begin{ujibox}
Setelah Langkah 16, periksa lima hal berikut.
\begin{itemize}[leftmargin=1.4em, itemsep=0.2em]
  \item Pemain berjalan mendatar, melompat hanya ketika menyentuh tanah, dan
        tidak dapat melompat berkali-kali di udara.
  \item Pemain dapat naik dari lantai ke rak pertama, lalu ke rak kedua.
  \item Artefak lenyap ketika disentuh, dan \menu{Console} mencatat jumlah
        artefak yang telah diambil.
  \item Pintu menolak pemain selama artefak belum lengkap, dan menerima
        pemain setelah ketiganya diambil.
  \item Lingkaran hijau pemeriksaan tanah tampak di bawah kaki pemain ketika
        \code{Pemain} dipilih pada tampilan \menu{Scene}.
\end{itemize}
\end{ujibox}

% ---------------------------------------------------------
\bagian{Kode Sesi Ini}

Tiga skrip berikut membagi pekerjaan dengan jelas. \code{KontrolPemain}
menangani masukan dan fisika pemain, \code{Artefak} menangani pengambilan
benda, dan \code{Pintu} menangani syarat keluar level.

\tampilkode[language=CSharp]{code/sesi-02/Scripts/KontrolPemain.cs}
  {KontrolPemain.cs: masukan pada Update, gerak dan pemeriksaan tanah pada FixedUpdate}

\tampilkode[language=CSharp]{code/sesi-02/Scripts/Artefak.cs}
  {Artefak.cs: trigger yang menandai artefak telah diambil}

\tampilkode[language=CSharp]{code/sesi-02/Scripts/Pintu.cs}
  {Pintu.cs: trigger yang memeriksa syarat artefak sebelum membuka level}

\subsection*{Tabel objek scene}

Tabel berikut adalah daftar objek yang harus ada pada scene Anda. Nilai yang
tidak disebut boleh dibiarkan bawaan.

{\footnotesize
\renewcommand{\arraystretch}{1.3}
\begin{tabularx}{\linewidth}{@{}L{2.4cm} L{3.5cm} Y@{}}
\toprule
\textbf{GameObject} & \textbf{Komponen} & \textbf{Nilai penting} \\
\midrule
Grid & Transform, Grid &
  Cell Size $(1,1,0)$ \\
TilemapLantai & Tilemap, Tilemap Renderer, Tilemap Collider 2D,
  Composite Collider 2D, Rigidbody2D &
  Layer \code{Ground}; Used By Composite; Body Type Static; petak
  $x=-9$--$9$ pada $y=-4$ \\
TilemapRak & Tilemap, Tilemap Renderer, Tilemap Collider 2D,
  Composite Collider 2D, Rigidbody2D &
  Layer \code{Ground}; Used By Composite; Body Type Static; tiga rak selebar
  empat petak \\
Pemain & Sprite Renderer, Rigidbody2D, Capsule Collider2D, KontrolPemain &
  Tag \code{Player}; Gravity Scale 2.5; Continuous; Freeze Rotation Z;
  Capsule $(0.7,0.9)$; Move Speed 6; Jump Speed 11; Ground Radius 0.15 \\
TitikTanah & Transform &
  Anak \code{Pemain}; local position $(0,-0.5,0)$ \\
Artefak1 & Sprite Renderer, Circle Collider2D, Artefak &
  $(-4,-1.1,0)$; $(0.5,0.5)$; \code{#DEB242}; Radius 0.25; Is Trigger;
  Nilai 1 \\
Artefak2 & Sprite Renderer, Circle Collider2D, Artefak &
  $(1,0.9,0)$; Radius 0.25; Is Trigger; Nilai 1 \\
Artefak3 & Sprite Renderer, Circle Collider2D, Artefak &
  $(6,-1.1,0)$; Radius 0.25; Is Trigger; Nilai 1 \\
Pintu & Sprite Renderer, Box Collider2D, Pintu &
  $(8.3,-2.6,0)$; $(1,1.8)$; Is Trigger; Artefak Diperlukan 3 \\
\bottomrule
\end{tabularx}
}

\checkpoint{02}{\raggedright Folder checkpoint memuat \file{Scripts/KontrolPemain.cs},
\file{Scripts/Artefak.cs}, \file{Scripts/Pintu.cs}, dan \file{scene.md}.
Scene tidak disertakan sebagai berkas biner; susunlah mengikuti tabel di atas
dan berkas \file{scene.md}.}

% ---------------------------------------------------------
\bagian{Kesalahan Umum}

\begin{galatbox}
\begin{itemize}[leftmargin=1.4em, itemsep=0.35em]
  \item \textbf{Pemain bergetar ketika berjalan.}
        Gejala: gambar pemain tampak gemetar, terutama ketika kamera
        mengikutinya. Periksa tempat \code{Rigidbody2D} digerakkan: bila
        melalui \code{transform.position} pada \code{Update}, pindahkan ke
        \code{body.velocity} pada \code{FixedUpdate}, lalu aktifkan
        \menu{Interpolate}.
  \item \textbf{Pemain menembus lantai atau rak.}
        Gejala: pemain jatuh melewati petak. Periksa \menu{Used By Composite}
        pada \code{Tilemap Collider 2D}, keberadaan \code{Composite Collider
        2D}, dan \menu{Collision Detection} pemain. Tampilkan bentuk batas
        melalui \menu{Window > Analysis > Physics Debug}.
  \item \textbf{Pemain jatuh tanpa henti.}
        Gejala: pemain melewati lantai dan terus turun. Periksa
        \menu{Ground Layer} pada \code{KontrolPemain} (harus \code{Ground}),
        \menu{Layer} kedua tilemap, dan \menu{Collision Matrix} pada
        \menu{Project Settings > Physics 2D}. Bila ketiganya benar, periksa
        apakah tilemap benar-benar memiliki collider.
  \item \textbf{Pemain dapat melompat berkali-kali di udara.}
        Gejala: tombol lompat yang ditekan berulang menghasilkan lompatan
        baru. Periksa apakah lompatan dibaca dengan
        \code{Input.GetButtonDown} dan apakah \code{onGround} diperiksa
        sebelum lompatan. Periksa pula \menu{Ground Radius}: lingkaran yang
        terlalu besar masih menyentuh rak di atas kepala.
  \item \textbf{Trigger tidak pernah dipanggil.}
        Gejala: artefak tidak lenyap dan Console tetap kosong. Periksa
        \menu{Is Trigger} pada collider artefak atau pintu, tag \code{Player}
        pada pemain, dan \menu{Layer Collision Matrix} antara layer pemain
        dan layer artefak.
  \item \textbf{Rujukan Inspector kosong.}
        Gejala: pesan \code{NullReferenceException} menunjuk baris
        \code{groundCheck.position}. Periksa \menu{Ground Check} pada
        komponen \code{KontrolPemain}; peubah itu harus berisi
        \code{TitikTanah}, bukan objek lain.
\end{itemize}
\end{galatbox}

% ---------------------------------------------------------
\bagian{Lembar Kerja}

Isi tabel berikut, lalu simpan sebagai bagian dari catatan praktikum Anda.

{\small
\renewcommand{\arraystretch}{1.4}
\begin{tabularx}{\linewidth}{@{}C{0.8cm} L{5.4cm} Y@{}}
\toprule
\textbf{No} & \textbf{Hasil yang dicatat} & \textbf{Keterangan Anda} \\
\midrule
1 & Nilai Gravity Scale dan Jump Speed yang dipakai & \\[0.55cm]
2 & Tinggi lompatan yang teramati, diukur dari lantai ke rak & \\[0.55cm]
3 & Jumlah petak lantai dan petak rak pada kedua tilemap & \\[0.55cm]
4 & Masalah fisika yang muncul beserta gejalanya & \\[0.9cm]
5 & Penyebab dan perbaikan masalah tersebut & \\[0.9cm]
6 & Pesan Console saat artefak diambil dan saat pintu dicoba & \\[0.9cm]
\bottomrule
\end{tabularx}
}

% ---------------------------------------------------------
\bagian{Checklist Sesi}

\begin{checklistbox}
\begin{checklist}
  \item Saya dapat menunjukkan lantai dan rak yang tersusun dari
        \istilah{Tilemap} dan memiliki satu bentuk batas utuh.
  \item Saya dapat menunjukkan pemain yang berjalan, melompat, dan mendarat
        di atas rak tanpa menembusnya.
  \item Saya dapat menjelaskan perbedaan \istilah{Collider} biasa dan
        \istilah{trigger}, beserta syarat agar trigger bekerja.
  \item Saya dapat menjelaskan mengapa masukan dibaca pada \code{Update}
        sedangkan gerak \istilah{Rigidbody2D} dikerjakan pada
        \code{FixedUpdate}.
  \item Saya dapat menunjukkan lingkaran pemeriksaan tanah dan menjelaskan
        fungsi \menu{Ground Layer}.
  \item Saya dapat menelusuri satu masalah fisika sampai penyebabnya, lalu
        memperbaikinya.
\end{checklist}
\end{checklistbox}

% ---------------------------------------------------------
\bagian{Pengait ke Sesi Berikutnya}

Level ini masih tampak sebagai balok dan petak berwarna. Pada sesi 3 Anda
mengimpor \istilah{sprite}, menyusun \istilah{Animator} untuk keadaan diam,
berlari, dan melompat, lalu membangun \istilah{Canvas} yang menampilkan skor
dan nyawa. Mekanik yang baru saja Anda selesaikan tidak diubah; yang berubah
hanya tampilannya.
